\section{Limitations}
\label{sec:limitations}

We acknowledge several limitations:

\paragraph{Scale and language.}
BiTempQA v3 contains 415 items (84 test) in Chinese only. While sufficient for diagnostic analysis, the small test set limits the granularity of per-type breakdowns (e.g., only 4 scenario types appear in the test set). Expanding the test set and adding English translations are priorities for future work.

\paragraph{Short-context scenarios only.}
Current scenarios contain 2--7 memory writes each, reflecting the length of ChronoQA golden chunks and curated event descriptions.
This limits our ability to test retrieval selectivity under longer contexts (50+ writes), where retrieval design may become a stronger differentiator.
The tight clustering of top systems (81.0--82.1\%) may not hold for longer scenarios.

\paragraph{Test set composition.}
The test set consists entirely of real-source scenarios and covers only 4 of 10 scenario types (entity attribute evolution, multi-source information, relationship evolution, temporal ambiguity).
The remaining 6 types are represented only in the training set via LLM-generated scenarios.

\paragraph{Abstractive-only evaluation.}
All test questions use abstractive (free-text) answers evaluated by LLM judge.
LLM judge leniency may inflate accuracy relative to human evaluation.

\paragraph{Production system LLM confound.}
Mem0 and Graphiti use Qwen2.5-32B-Instruct as their internal LLM for memory extraction, while lightweight backends use DeepSeek-V3 for answer generation.
Their low scores (2--8\%) therefore reflect a combination of architectural limitations (memory compression, lack of bitemporal indexing) and LLM integration challenges.
Evaluating with a fully compatible LLM backend may yield higher absolute scores, though we expect the qualitative patterns---compressed memories losing temporal detail---to persist.

\paragraph{ChromaDB interpretation.}
While ChromaDB uses the same embedding model as other vector systems, its metadata filtering approach is architecturally distinct from pure similarity search. Its 22.6\% accuracy reflects the interaction between metadata filtering and heterogeneous temporal formats, not necessarily a fundamental limitation of vector databases.
